\scanpage{180}\mspage{35}
n'y a plus lieu de retenir que le ``type d'homotopie''. On prévoit que les 6 opérations sur ces faisceaux peuvent se traduire ad-hoc dans ce cadre homotopique. Enfin, si au lieu d'un seul espace modéré compact $X$, on a disons un morphisme modéré $f:X\unsure{-|\to}Y$ de tels, alors, choisissant des stratifications équisingulières sur $X$ et $Y$ adaptées à $f$ (les strates de $X$ étant notamment des fibrés loc. triviaux sur celles de $Y\ldots$), on doit trouver un ``morphisme'' du diagramme de variétés à bord exprimant $X$ dans celui exprimant $Y$ (avec des morphismes verticaux se ramenant \deleted{\unsure{à|,}} encore, pour l'essentiel, à des fibrations de variétés à bord sur d'autres) de telle façon que les 4 opérations $\mathrm{R}\deleted{\unsure{f}}g_*$, $\mathrm{R}g_!$, $\mathrm{L}g^*$, $g^!$ entre faisceaux \inserted{$\unsure{P|\mathcal{P}}_X$- et $P_Y$-constructibles} sur $X$ et $Y$ (ou sur des sous-espaces loc. fermés \inserted{$X'$, $Y'$} réunions de strates, de sorte que $f$\unsure{,|.} induise $g:X'\unsure{-|\to}Y'$) s'expriment encore en termes d'opérations standart entre \unsure{seuls|de tels} types d'homotopie. Enfin, toutes ces constructions, encore partiellement hypothétiques (y a du travail \inserted{de fondements} sur la planche!) doivent
% archive sheet number: 180
